\begin{proof}
(i) Premise E (\claimref{dark-matter.e-equal-mass-per-unit-ball}) for a three-dimensional tower and for this world reads $\Sigma\,V_{2}d^{2} = \rho_{b}V_{3}d^{3}$, with $V_{2} = \pi$ and $V_{3} = \tfrac{4}{3}\pi$, so $\Sigma = \tfrac{4}{3}\rho_{b}d$; the cells do not appear. With the sheets $d$ apart (Premise R3, \claimref{dark-matter.r3-one-cell-per-unshared-direction}) and comoving, one per unit of comoving length (Lemma \claimref{dark-matter.projection-every-towers-content-dilutes-ones}), $\Sigma$ is a comoving areal density. At $d = d_{b}$, $\rho_{b} = 3m_{N}/(4\pi d_{b}^{3})$ gives $\Sigma\,\pi d_{b}^{2} = m_{N}$.

(ii) On a sheet the number per unit area is $n = N a^{-2}$ with $N$ comoving, and in radiation domination $H = h\,a^{-2}$ with $h = a^{2}H$ constant at fixed $g$. The rate equation $\dot n = -2Hn - \sigma_{1}v\,n^{2}$ becomes $dN/d\ln a = -(\sigma_{1}v/h)\,N^{2}$, whose solution is $1/N = 1/N_{0} + (\sigma_{1}v/h)\ln(a/a_{i})$. With $\sigma_{1}v = 2P/m$ and $N_{0} = \Sigma/m$ the depletion term $N_{0}/N_{f} - 1$ is $2P\,\Sigma\ln(a_{f}/a_{i})/(m^{2}h)$. By (i) $\Sigma$ is proportional to $d$, so the term is proportional to $P\,d/m^{2}$, that is to $P X$. The bound $P_{\max}$ is the value of $P$ at which the term is $1/99$ at $X = 1$, evaluated in rational intervals in the check.

(iii) Observation \claimref{dark-tower-matter.three-dimensional-phases} gives contact processes the scaling $(m/e^{2})^{2}$ in phase (T) for $e^{2} \gg m$. Substituting $P = (m/e^{2})^{2}$ in (ii), $m^{2}$ cancels and the term is $2\Sigma\ln(a_{f}/a_{i})/(e^{4}h)$. It is at most $1/99$ exactly when $e^{4} \ge 198\,\Sigma\ln(a_{f}/a_{i})/h$, and $\Sigma = (m_{N}/\pi d_{b}^{2})(d/d_{b})$ puts this in the form $e^{2} \ge c\,m_{N}\sqrt{d/d_{b}}$, $c = P_{\max}^{-1/2}$. At $m = m_{N}$, $d = d_{b}$ the same inequality reads $m/e^{2} \le P_{\max}^{1/2}$.

(iv) A cell collides only within its sheet, at the rate $n_{2}\lambda v = (n_{2}/s)(\lambda s)v$, the rate of a gas of number density $n_{2}/s$ and cross-section $\lambda s$; the mass density that gravity sees is $n_{2}m/s$ by the sourcing rule (Lemma \claimref{dark-matter.projection-every-towers-content-dilutes-ones}). By Fact \claimref{simulations.imported-low-energy-scattering-two-dimensions}, $\lambda = (4/k)\sin^{2}\delta_{0}$ at low energy, with $k = mv/2$ for equal masses, so $\sigma_{\mathrm{eff}}/m = 8\sin^{2}\delta_{0}\,s/(m^{2}v)$. The spacing in a region of given compression is proportional to $d$, so this is proportional to $\sin^{2}\delta_{0}\,d/m^{2}$. The bound is that of Fact \claimref{simulations.imported-bullet-cluster}, taken at $1\ \mathrm{cm^{2}\,g^{-1}}$ as T14b takes it; $S_{*}$ is evaluated in rational intervals in the check. With $\sin\delta_{0} = m/e^{2}$, $m^{2}$ cancels as in (iii).

(v) Setting $P = 1$ in (ii) gives $X \le P_{\max}$, that is $m \ge c\,m_{N}\sqrt{d/d_{b}}$; higher partial waves only add to the rate, so it is a lower bound. For clusters, the low-energy form of the cited Fact has $\lambda \approx \pi^{2}/(k\ln^{2}(ka))$, a function of $ka$ over $k$: at fixed $ka$ it goes as $1/m$ and $\sigma_{\mathrm{eff}}/m$ as $d/m^{2}$; for $ka \gg 1$ a disc's cross-section is geometric, of order $a$, and $\sigma_{\mathrm{eff}}/m$ goes as $a\,d/m$.

(vi) E states one length common to all populated towers, and R3 takes it as the thickness of both kinds of odd tower. In the stated model the four-dimensional coupling at depth $L$ is $g_{5}^{2}/L$ and the one-loop exponent is $8\pi^{2}L/(11\,g_{5}^{2}) = m_{5}L/(11\pi)$; with $L$ comoving, $L = 20\,(LT)/m_{5}$ at $T = m_{5}/20$, so the exponent is $20\,(LT)/(11\pi)$, in which $m_{5}$ appears only through the number of species at that temperature, $LT = (d\,T_{0})(3.91/g_{s})^{1/3}$. The intervals for $\kappa$ and the two comparisons with the Hubble radius are evaluated in the check.
\end{proof}
